% ============================================================================
% Chapter 5: Implementation and Experimental Setup
% ============================================================================
\chapter{Implementation and Experimental Setup}
\label{ch:implementation}

\section{Technology Stack}

CBSR is implemented in Python 3.9 using the following core libraries:
\begin{itemize}
\item \textbf{Tigramite 5.x}: PCMCI causal discovery with RobustParCorr (MCD-based partial correlation).
\item \textbf{NumPy / SciPy}: Numerical computation, statistical functions, and CDF evaluations.
\item \textbf{PyTorch 2.0}: LSTM Autoencoder for Temporal State Reasoning ($R_T$).
\item \textbf{Scikit-Learn 1.3}: Behavioral Pattern Reasoner (Random Forest) and stacked meta-learner (Logistic Regression).
\item \textbf{Pandas}: Data ingestion and provenance edge preprocessing.
\end{itemize}

The complete implementation resides in two primary source files:
\texttt{zerocausal\_core.py} ($\approx$62\,KB) — core module containing all algorithms; and
\texttt{15\_beat\_papers\_evaluation.py} ($\approx$28\,KB) — end-to-end evaluation pipeline.


\section{Benchmark Datasets}
\label{sec:datasets}

CBSR is evaluated on five diverse provenance-based IDS benchmarks spanning simulated and
real-world enterprise environments. Table~\ref{tab:datasets} summarizes the key statistics.

\begin{table}[htbp]
\centering
\small
\caption{Benchmark dataset statistics. Windows are non-overlapping time slices of width $W$.
Train/Val/Test follow a strict 50\%/20\%/30\% chronological split.}
\begin{tabularx}{\linewidth}{Xrrrcrrr}
\toprule
\textbf{Dataset} & \textbf{Total} & \textbf{Train} & \textbf{Val} & \textbf{Test}
                 & \textbf{Attacks} & $d$ & $W$ \\
\midrule
DARPA TC3      & 399   & 200 & 80  & 119 & 134 & 23  & 1\,s \\
NODLINK        & 399   & 200 & 80  & 119 & 119 & 23  & 1\,s \\
BETH (Sysdig)  & 539   & 270 & 108 & 161 &   2 & 882 & 10\,s \\
DARPA OpTC     & 546   & 273 & 109 & 164 &  50 & 130 & 1\,s \\
StreamSpot     & 37{,}614 & 18{,}807 & 7{,}523 & 11{,}284 & 2{,}843 & 82 & --- \\
\bottomrule
\end{tabularx}
\label{tab:datasets}
\end{table}

\begin{itemize}
\item \textbf{DARPA TC3} (DARPA ENGAGE): A simulated Windows enterprise APT campaign with
  multi-stage attack chains --- persistence, lateral movement, credential dumping, and
  exfiltration. The provenance graph covers 23 unique edge types (process--file,
  process--network, process--process). AUC target: 0.9837.

\item \textbf{NODLINK} (DARPA NODLINK): A simulated campus-network APT trace captured via
  Linux audit logging. Attack stages include initial compromise, privilege escalation, and
  data exfiltration. Identical dimensionality to TC3 ($d = 23$) due to shared simulation
  infrastructure. AUC target: 0.9784.

\item \textbf{BETH} (USENIX ATC'22): Real-world Linux host telemetry captured via sysdig
  across 8 physical hosts, of which 3 experienced active red-team attacks. The dataset has
  $d = 882$ unique syscall-derived edge types and only 2 labeled attack windows in the test
  partition, making it an extremely imbalanced detection challenge. AUC target: 0.9901.

\item \textbf{DARPA OpTC}: Real DARPA enterprise APT evaluation across a live Windows
  environment of 500+ hosts. The provenance graph captures 130 unique edge types.
  50 attack-labeled windows are present in the full evaluation trace. AUC target: 0.9628.

\item \textbf{StreamSpot}: A graph-stream benchmark containing 600 heterogeneous provenance
  graphs across 6 distinct workload types (YouTube, GMail, download, VGame, system, attack).
  The attack class (100 graphs) is embedded in the StreamSpot-B split. With 37,614
  temporal windows at $d = 82$ edge types, it is the largest benchmark and the hardest due
  to workload heterogeneity. AUC target: 0.6909.
\end{itemize}


\section{Experimental Protocol}
\label{sec:protocol}

\subsection{Hardware and Software Environment}

\begin{table}[htbp]
\centering
\small
\caption{Hardware and software environment.}
\begin{tabular}{ll}
\toprule
\textbf{Component} & \textbf{Specification} \\
\midrule
Processor    & Intel Xeon E5-2698 v4, 20 cores @ 2.2\,GHz \\
Memory       & 128\,GB DDR4 ECC \\
Storage      & 2\,TB NVMe SSD \\
GPU          & NVIDIA Tesla V100 (32\,GB) — used for LSTM-AE training only \\
OS           & Ubuntu 20.04 LTS (Linux kernel 5.4.0) \\
Python       & 3.9.18 \\
CUDA         & 11.8 (PyTorch 2.0.1+cu118) \\
\bottomrule
\end{tabular}
\label{tab:hardware}
\end{table}

\subsection{Chronological Data Partitioning}

All experiments use a strict chronological 50\%/20\%/30\% three-way split to prevent
future-leakage, as required by any realistic streaming IDS evaluation:

\begin{enumerate}
\item \textbf{Train Proper (50\%)}: Used exclusively for unsupervised RobustParCorr causal
  discovery, SCM regression fitting ($\beta_{i,j}$), and LSTM Autoencoder pre-training on
  normal-period state sequences. No labels or attack-free guarantees are assumed.

\item \textbf{Validation (50\%--70\%)}: Provides the labeled partition on which the
  Behavioral Pattern Reasoner (BPR) and stacked meta-learner are supervised. Conformal
  calibration scores are also initialized here.

\item \textbf{Test (final 30\%)}: Strictly held out. Streaming risk inference with online
  conformal alerting is performed window-by-window; final AUC-ROC and FPR@95\%TPR metrics
  are computed on this partition only.
\end{enumerate}

\subsection{Reproducibility and Confidence Estimation}

All experiments are repeated across \textbf{5 random seeds} (42--46) that control LSTM-AE
initialization and RF training stochasticity. Mean AUC and 95\% confidence intervals
($\pm$1.96$\times$std) are reported. PCMCI causal discovery is cached per
\texttt{(dataset, seed, $\alpha$, $\tau_{\max}$)} to eliminate redundant computation across
evaluation runs.

\subsection{Evaluation Metrics}

\begin{itemize}
\item \textbf{AUC-ROC}: The primary metric; measures discrimination across all thresholds.
  Preferred over accuracy due to high class imbalance in all five benchmarks.
\item \textbf{FPR@95\%TPR}: The false positive rate when the detection threshold is set to
  achieve 95\% recall. This is the primary operational metric for a deployed IDS, where missing
  an attack (low TPR) is unacceptable.
\item \textbf{Conformal FPR}: Empirical false alarm rate under online conformal alerting at
  target $\alpha = 5\%$, reported in the noise and drift robustness analyses.
\end{itemize}


\section{Operator Architecture Details}

\textbf{Temporal State Reasoning (LSTM-AE):} Encoder LSTM (12-dim $\to$ 32-dim hidden) and
Decoder LSTM (32-dim $\to$ 12-dim output), trained for 60 epochs with Adam ($\text{lr} =
5 \times 10^{-4}$, batch size 64) on sequences of $k = 8$ consecutive normal-period state
vectors from the Train Proper partition.

\textbf{Behavioral Pattern Reasoner (BPR):} Random Forest with 300 decision trees, maximum
depth 10, and Gini impurity criterion, trained on the 12-dimensional state coordinates from
the Validation partition with corresponding binary attack labels.

\textbf{Calibrated Meta-Learner:} Logistic Regression with L2 regularization ($C = 1.0$),
fitted on the Validation partition's three reasoning operator outputs ($R_L$, $R_T$, $R_B$)
and labels. A rolling queue of $M = 245$ calibration scores is maintained online for
conformal p-value computation.


\section{Performance Optimizations}

When initially profiled, the streaming loop required 240\,s to process 546 windows on OpTC
($\approx$440\,ms/window). Three targeted optimizations achieved a \textbf{13$\times$ speedup}:

\textbf{NumPy Matrix Indexing.} Pandas \texttt{.iloc} calls inside the sliding window loop
added $\approx$200\,$\mu$s/access. Solution: pre-convert dataframes to 2D NumPy matrices and
map column names to integer offsets:
\begin{lstlisting}
data_arr = ts.values.astype(float)       # (N, d) matrix
col_idx = {c: i for i, c in enumerate(ts.columns)}
\end{lstlisting}

\textbf{Binary Search Conformal Lookup.} The original $O(M)$ scan of $M = 245$ calibration
scores was replaced with a sorted array and $O(\log M)$ binary search:
\begin{lstlisting}
calib_sorted = np.sort(calib_array)
rank = np.searchsorted(calib_sorted, score, side='right')
pval = (M - rank + 1) / (M + 1)
\end{lstlisting}

\textbf{Vectorized Empirical CDF.} Per-step SciPy CDF evaluations replaced with batch NumPy
operations, eliminating Python interpreter overhead inside the hot loop.


\section{Inference Latency}

These optimizations reduce total processing time from 240\,s to 18.03\,s on OpTC
($<$33.4\,ms/window for causal scoring). The full CBSR hybrid ensemble (three reasoning
operators + meta-fusion + conformal lookup) then runs at \textbf{1.35--1.71\,ms/window}
across all benchmarks (Table~\ref{tab:runtime}) --- well below the 1-second streaming window.

\begin{table}[htbp]
\centering
\caption{CBSR hybrid ensemble inference latency per streaming window.}
\begin{tabular}{lc}
\toprule
\textbf{Dataset} & \textbf{Avg Inference (ms)} \\
\midrule
TC3        & 1.705 \\
NODLINK    & 1.366 \\
BETH       & 1.351 \\
OpTC       & 1.354 \\
StreamSpot & 1.703 \\
\bottomrule
\end{tabular}
\label{tab:runtime}
\end{table}

PCMCI discovery scales as $O(d^2 \tau_{\max} T)$ in practice. On BETH ($d = 882$, $T = 270$
training windows), the one-time offline discovery completes in $\approx$7,842\,s and is cached
via \texttt{.pkl} for all subsequent seed runs. For $d > 500$, a variance-thresholding pre-pass
reduces effective dimensionality to maintain sub-hour startup. The figure below plots discovery
time (log scale) and per-window latency as functions of $d$: discovery grows as $O(d^2)$ while
inference latency stays constant at $\approx$1.4\,ms regardless of $d$.

\begin{center}
\includegraphics[width=0.56\linewidth]{dimensionality_vs_latency.png}
\captionof{figure}{PCMCI discovery time and streaming inference latency vs.\ $d$.}
\label{fig:scalability}
\end{center}
